\SourceLeaf{417}%
in der That für den in Rede stehenden singulären Fall immer vier Auflösungen, indem man entweder durch \((4)\) oder durch \((5)\) eine der beiden geraden Linien legt, deren Winkel mit der Abscissen-Linie die Grösse \(\frac A{2a}=\frac B{2b}\) zum Cosinus hat. Die Realität dieser Auflösungen hängt davon ab, dass diese Grösse nicht grösser als \(1\) ist, für welchen Werth selbst die vier Auflösungen sich auf zwei reduciren.
\par\bigskip\begin{center}\rule{.16\textwidth}{.3pt}\end{center}
